% !TEX root = ../main.tex
% !TEX program = lualatex
\documentclass[../main.tex]{subfiles}
\IfSubfilesClassLoaded{\externaldocument{\subfix{../build/main}}}{}
% =============================================================================
% APPENDIX: CNO 34 DAY ONE MESSAGE
% DESCRIPTION: Full text of ADM Caudle's Foundry--Fleet--Fight day-one message.
% =============================================================================

\begin{document}
\IfSubfilesClassLoaded{\chapter{EDO Study Guide}}{}

\chapter{CNO Strategic Messaging}\label{ch:cno-strat}
\section{CNO 34 Day One Message (Foundry--Fleet--Fight)}\label{app:cno34-day-one-message}
\emph{\footnotesize NAVADMIN 180/25; issued upon assuming duties as the 34th \ac{cno}~\autocite{cno34-day-one-message-2025}.}

\begin{description}[style=nextline, labelsep=0.5em, font=\bfseries]
  \item[\textbf{Mission.}] The Navy shall be organized, trained, and equipped to deliver peace through strength, securing our national interests and prosperity. We exist for prompt and sustained combat at sea, preserving freedom of navigation, safeguarding global commerce, and deterring those who would challenge America's values and sovereignty.

  \item[\textbf{Vision.}] We stand at an inflection point---an era marked by great power competition, proliferating threats, rapid technological convergence, and an increasingly contested maritime domain. To prevail, we must build and sustain a Navy that is ready to fight and win---today, tomorrow, and well into the future. At the center of this vision is the Sailor. Our men and women are the Navy's main weapon system---their skill, resilience, and courage remain our greatest asymmetric advantage. We will put them first by relentlessly pursuing full-spectrum readiness, deepening cross-domain integration, harnessing innovation at speed, and strengthening our warrior ethos. Our Navy must be resilient, agile, globally present, and combat credible. We will align tightly with national defense priorities, expand our ability to integrate with Joint Force and Allied partners, and preserve our enduring advantage at sea. In doing so, we will guarantee freedom of navigation, deter aggression, and project power where and when it matters most.

  \item[\textbf{Priorities.}] We will view everything we do through an operational lens focused on the Foundry, the Fleet, and the way we Fight.
        \begin{description}[style=nextline, labelsep=0.5em, font=\bfseries]
          \item[\textbf{Foundry.}] The Foundry represents the enduring sum of our total force, shore infrastructure, maintenance depots, schoolhouses, industrial base, and intellectual capital required to generate, sustain, and modernize naval power. Every ounce of sea power we wield originates in the Foundry; it is the bedrock of our Navy. We will not neglect the global network of infrastructure and capacity that enables our Sailors to excel. From new ship construction to depot-level repair, from weapons production to the training pipelines that develop warfighters, the Foundry underwrites combat credibility. To remain the world's most formidable fighting force, we will operationalize readiness, accelerate shipbuilding and repair at scale, and modernize our force-generation model. This means strengthening public-private partnerships, investing in resilient supply chains, expanding workforce development, and ensuring our training systems deliver Sailors fully prepared for the high-end fight. Simply put---without a strong Foundry, there is no Fleet.

          \item[\textbf{Fleet.}] The Fleet is forged in the Foundry and is our Nation's decisive instrument of maritime power. Comprised of people, platforms, and payloads, it is designed to deter aggression, project power, and secure our national interests on and from the sea. No image symbolizes American resolve more than a U.S. Navy ship on the horizon---lethal, resilient, and ready to fight. The Fleet's unique mobility, persistence, and global reach enable us to impose costs on adversaries, reassure allies, and guarantee the free flow of commerce. Tempered by our resolve, our Fleet is world-class, battle-ready, and second to none. With our Sailors at the helm, supported by strong Navy families, there is no challenge we cannot overcome. We will continue to build an integrated, all-domain Fleet---one that fuses maritime (surface, air, and undersea), cyber, space, and strategic capabilities while leveraging Joint Force and Allied partners to dominate across the competition continuum. Interoperability and interchangeability with allies and partners will be our force multiplier. By enhancing collective maritime security, we will magnify our forward presence and strengthen deterrence worldwide.

          \item[\textbf{Fight.}] Fighting is our ultimate honor and sacred responsibility---delivering violence on behalf of the Nation when called. For nearly 250 years, the U.S. Navy has honed its toughness, tenacity, and ingenuity to secure victory. Those who challenge us in combat will be met with overwhelming force and sent to the bottom. But we must fight differently in the future. In today's rapidly evolving landscape, we will innovate, adapt, and integrate cutting-edge technologies at the speed of relevance. We will deliver and refine the Navy Warfighting Concept and develop a new Navy Deterrence Concept, both anchored in peace through strength. We will accelerate the integration of artificial intelligence, robotic and autonomous systems, resilient \ac{c3} architectures, and synchronized effects---ensuring decision advantage, survivability, and combat overmatch. We will ruthlessly pursue a Future Fleet Design---comprised of unmatched conventionally manned ships, deeply integrated with asymmetric hedge forces that expand our reach and impose costs more effectively across every domain. Our charge is clear: To ensure the Navy remains the most lethal, survivable, and globally dominant maritime force.
        \end{description}

  \item[\textbf{Theory of victory.}] At the end of my tenure as \ac{cno}, success will be measured by results:
        \begin{itemize}
          \item Platforms delivered and repaired on time.
          \item Fully manned and combat-ready ships.
          \item Ordnance production meeting contracted demand.
          \item Backlogs in repair parts eliminated.
          \item Sailors trained to the highest levels of mastery.
        \end{itemize}
        Through disciplined execution, we will increase our lethality and preserve our sustained advantage. To achieve this, we are calibrating our sights, refining our course, and transforming our approach of solving our Navy's toughest problems. Guided by the Navy Warfighting and Deterrence Concepts, we will develop and execute a holistic Campaign Plan that delivers the warfighting capability and capacity necessary to win in the 21st century. To maintain transparency and unity of effort, I will release a series of ``CNO Notes''---or ``C-Notes''---providing guidance and direction to the Fleet. These will be grounded in my vision, informed by your feedback, and driven by operational necessity. They will require your trust, ownership, and commitment to be successful. I am proud to serve alongside you---America's Sailors---as we embark on this next chapter in our Navy's history. Built in the Foundry---Tempered in the Fleet---Forged to Fight.

        \medskip
        Daryl L.\ Caudle\\
        Admiral, United States Navy\\
        34th Chief of Naval Operations
\end{description}

\section{CNO 34 C-Notes and Charge of Command}\label{app:cno34-cnotes}
\emph{\footnotesize NAVADMINs 186/25, 203/25, 241/25, and 225/25~\autocite{cno-cnote-1-2025,cno-cnote-2-2025,cno-cnote-3-2025,cno-charge-of-command-2025}.}

\subsection{Executive summary of priorities}
\begin{description}[style=nextline, labelsep=0.5em, font=\bfseries]
  \item[Overview.] ADM Caudle's initial set of C-Notes and his Charge of Command reinforce the Foundry--Fleet--Fight construct: care for Sailors, strengthen the Navy--industrial team, and deliver a combat-credible fleet that can deter and win. Leaders are directed to build trust, enforce standards, and maintain a visible warfighting mindset across every command.
  \item[Overview.] \ac{adm} Caudle's initial set of C-Notes and his Charge of Command reinforce the Foundry--Fleet--Fight construct: care for Sailors, strengthen the Navy--industrial team, and deliver a combat-credible fleet that can deter and win. Leaders are directed to build trust, enforce standards, and maintain a visible warfighting mindset across every command.
  \item[C-Note \#1: Sailors First.] The Fleet's advantage is a trained, supported Total Force. Priorities include modern housing and galley upgrades, streamlined uniform requirements, resilient digital connectivity ashore, and pay/benefits accuracy so every Sailor can focus on the fight~\autocite{cno-cnote-1-2025}.
  \item[C-Note \#2: Foundry Always.] The Foundry---people, infrastructure, and materiel---forges naval power. Caudle calls for relentless training (Career Training Continuum), expanded live-virtual-constructive integration, stronger shipyard shifts, revitalized intermediate maintenance activities, and modernized shore platforms and supply chains to deliver ready forces at scale~\autocite{cno-cnote-2-2025}.
  \item[C-Note \#3: World Class Fleet.] The Fleet is the decisive instrument of national power, integrating people, platforms, and payloads across all domains. The note introduces a Hedge Strategy (tailored offsets and tailored forces) to outpace threats, modular payload investments, resilient \ac{c2}/\ac{ai}/unmanned integration, clarified administrative/operational control, and undersea \ac{c2} realignments (\ac{ctf}-54/74)~\autocite{cno-cnote-3-2025}.
  \item[Charge of Command.] Command is an action and a privilege. Commanding Officers are accountable for warfighting readiness, disciplined risk management, mission command, continuous learning, and visible ownership of their Sailors' welfare and performance. Immediate superiors will review the Charge with subordinate \acp{co} to ensure shared understanding and mentorship~\autocite{cno-charge-of-command-2025}.
\end{description}

\subsection{Message texts}
\paragraph{C-Note \#1: Sailors First (NAVADMIN 186/25).}
\begin{quote}
  \small\begin{verbatim}
CLASSIFICATION: UNCLASSIFIED/
ROUTINE
R 031453Z SEP 25 MID180002011151U
FM CNO WASHINGTON DC
TO NAVADMIN
INFO CNO WASHINGTON DC
BT
UNCLAS
NAVADMIN 186/25
MSGID/NAVADMIN/CNO WASHINGTON DC/CNO/SEP//
SUBJ/C-Note #1: Sailors First//
RMKS/1. The lifeblood of our Navy is the Sailor: well-trained, connected,
supported and fit to fight. In every ocean around the world, our national
security interests are challenged by increasingly complex threats.
Underpinned by our warrior ethos, we will rise to that challenge. For 250
years, we have been a world-class fighting Navy, and our strength lies in our
Sailors. Our teamwork and esprit de corps are strengthened by our core
values and devotion to our mission.
2. We must resource the Navy our Nation Needs. I am committed to the well-
being of our Sailors and their families - providing them with state-of-the-
art platforms, world-class facilities and dependable support through our
Total Sailor and Family Action Plan.
We will leverage the insights and lessons learned from the past to improve
and rebrand our Culture of Excellence initiative to its new focus - Total
Sailor: Fit to Fight.
3. Today, I have directed your Navy leadership to take the steps necessary
to ensure that No Sailor Will Live Afloat. We will invest in unaccompanied
and family housing in addition to optimizing our new BAH authority to ensure
every Sailor resides in housing that is clean, comfortable and safe.
4. During my term, we will transform our galley facilities to ensure they
are designed to meet the nutritional needs and choices of our Sailor's
modern-day lifestyle. It is paramount that we modernize food service to
deliver healthy and flexible options.
5. Effective immediately, the Secretary of the Navy has directed we convene
a focused assessment of your uniforms and seabag. We will streamline uniform
requirements to ensure you have affordable, available and durable service
clothing that you are proud to wear.
The goal will be to reduce the number of uniforms while respecting our hard-
earned traditions that makes our Navy so special.
6. You have grown up in a digitally connected world where disconnection
means isolation. Therefore, I have mandated efforts to improve cell service
on our installations and bring free Wi-Fi to all of our barracks.
7. I have also ordered changes to the NAVADMIN process ensuring our Sailors
are able to find clear and correct information in one singular location. We
will not use NAVADMINs to work around the effort required to properly change
the base instructions any longer.
8. Additionally, I've directed Navy leadership to revisit prior efforts such
as Sailor pay to ensure our previous corrective actions are still working so
that each Sailor is receiving the correct amount of pay and entitlements in a
timely manner.
9. You are my priority and I need you focused on the fight. This is the
first of many efforts to improve your Quality of Life - so, stand by for
what's coming next.
10. Built in the Foundry - Tempered in the Fleet - Forged to Fight.
11. ADM Daryl Caudle, 34th Chief of Naval Operations sends.
BT
#0001
NNNN
CLASSIFICATION: UNCLASSIFIED/
\end{verbatim}
\end{quote}

\paragraph{C-Note \#2: Foundry Always (NAVADMIN 203/25).}
\begin{quote}
  \small\begin{verbatim}
CLASSIFICATION: UNCLASSIFIED/
ROUTINE
R 222218Z SEP 25 MID120012080004U
FM CNO WASHINGTON DC
TO NAVADMIN
INFO CNO WASHINGTON DC
BT
UNCLAS
NAVADMIN 203/25
MSGID/NAVADMIN/CNO WASHINGTON DC/CNO/SEP//
SUBJ/C-NOte #2: Foundry Always//
RMKS/1. The Foundry is the engine that drives our warfighting advantage. It
transforms raw inputs and forges them into lethal
outputs: ready warfighters, platforms and sustained readiness, at the speed
and scale necessary to win.
2. The Foundry builds, generates, sustains and modernizes naval power
through three critical components: PEOPLE, INFRASTRUCTURE and MATERIEL.
3. People: Our Total Force, comprised of Sailors, civilians and the
industrial workforce, is our most decisive warfighting advantage.
To prevail in conflict, we must recruit, train, educate and retain a
technically proficient and resilient team. Investing in our people's
development and well-being is not separate from combat power, it is the very
foundation of our Warrior Ethos, which is not exclusive to those in uniform.
It's a mindset marked by honor, integrity, resilience and a relentless
commitment to excellence and service.
That ethos must live across the spectrum: from our schoolhouses and combat
centers to the workforce that lays our keels and develops our intellectual
property. Every role our Total Force fills is vital by design.
4. To keep pace with the proliferation of technology and the global
evolution of threats, we are enhancing the quality of our curricula.
Relentless improvement, technical competence and tactical mastery are all
tenets demanded of a confident warfighter. Empowered by timely action and
responses from our Navy's top leaders, we will inculcate a culture of
continuous learning that delivers world-class, modernized and battle-ready
Sailors. Therefore, we will rename Ready Relevant Learning to "Career
Training Continuum" (CTC), the Navy's enterprise program of record for
individual training and career-long learning.
5. Through a series of sprints, we will also expand the Military Learning
Continuum and Civilian Workforce Development. With built-in feedback loops
and mechanisms in place, we will rapidly identify, assess, respond to and
correct gaps in our Sailor's training journey.
I have also directed that our Live-Virtual-Constructive (LVC) efforts be
expanded past the current use cases (e.g., integrated C2X training, Fallon
airwing events, etc.) to include a more generalized view of LVC that spans
from unit level training to schoolhouse utilization to tailored wargaming.
6. Battle effectiveness is not achieved by only afloat platforms or
commands. Simply put, without a strong Foundry, there is no effective Fleet.
Therefore, I am working to rewrite the Battle "E"
instruction to recognize shore installations with exceptional readiness and
performance.
7. Infrastructure: Shore platforms generate naval power. Our global
network of infrastructure, shipyards, airfields, barracks, electrical grids
and piers are foundational to projecting that power. Modernization ensures
our readiness can scale across the spectrum of conflict. We will prioritize
and strengthen our shore platforms to support Fleet-wide operations and
ensure our People take pride and ownership in the places that they work in
and live.
8. Expanding the Fleet's deployable platforms is a whole-of-Navy evolution,
and while we are not short on ideas to get there, our execution has not met
expectations. The levers: innovation, oversight, optimization and
recapitalization are known, but the fastest gains in production will come
from our People and their elbow-grease. To meet operational timelines, we
must prioritize workforce development, labor capacity and world-class project
management to increase our efficiency.
9. Accordingly, we will strengthen second and third shifts in our public
shipyards to increase throughput and on-time delivery. Only by improving
shipyard training, workplace conditions and organizational support can we
reduce attrition, accelerate experience and retain the workforce, both
private and public, responsible for delivering our Fleet.
10. I will release a future C-NOte on Esprit de Corps Ashore, emphasizing
the critical role our infrastructure plays in forging warfighting readiness.
11. Materiel: Reliable supply chains, accurate Coordinated Shipboard
Allowance Lists (COSALs), comprehensive maintenance capabilities and
coordinated logistics are instrumental to keeping our Fleet in the fight.
Winning in conflict demands a strong Navy-Industrial Base partnership,
grounded in transparency, common understanding and intellectual honesty.
12. Our current maintenance structure lacks the right balance of Sailors,
civilians, artisans and engineers. Years of depot and activity level
consolidation have reduced opportunities for our Sailors to master the skill
sets needed to repair and sustain at sea.
13. For these reasons, I am pursuing the stand-up of Shore Intermediate
Maintenance Activities (SIMA) in Norfolk and San Diego.
These facilities will provide our Sailors with hands-on training in advanced
ship repair and expose them to modern capabilities such as AI/ML, advanced
manufacturing (e.g., 3-D printing, CAD/CAM, etc.), workflow monitoring
technologies and robotic systems.
14. We are also reviewing Class Maintenance Plans. To maintain a combat
surge-ready posture, our maintenance models must evolve.
Introducing shorter, more frequent and efficient availabilities will
increase our readiness, improve on-time completion and reduce risk at a
comparable cost.
15. Foundry is more than a metaphor; it is the process, standard and
discipline that underpins our Navy's ability to deliver ready, lethal and
resilient forces, on time, on cost and at scale.
16. Built in the Foundry - Tempered in the Fleet - Forged to Fight.
17. ADM Daryl Caudle, 34th Chief of Naval Operations sends.//
BT
#0001
NNNN
CLASSIFICATION: UNCLASSIFIED/
\end{verbatim}
\end{quote}

\paragraph{C-Note \#3: World Class Fleet (NAVADMIN 241/25).}
\begin{quote}
  \small\begin{verbatim}
CLASSIFICATION: UNCLASSIFIED/
ROUTINE
R 012011Z DEC 25 MID120012255382U
FM CNO WASHINGTON DC
TO NAVADMIN
INFO CNO WASHINGTON DC
BT
UNCLAS
NAVADMIN 241/25
MSGID/NAVADMIN/CNO WASHINGTON DC/CNO/DEC//
SUBJ/C-Note #3: World-Class Fleet//
RMKS/1. History proves that great nations are built on the strength of great
navies. And today, our Fleet stands as the world's premier warfighting force
- ready, capable, and unmatched. But we are more than ships, aircraft, and
weapons. Our Sailors - battle ready, trained, resilient, and forged in the
Foundry are the engine that empowers our Fleet. Through their strength, we
command the seas, deter aggression, defend our homeland, and when called,
deliver decisive combat power. This is what makes us a world-class Navy -
second to none.
2. The Fleet is our most decisive instrument of national power and serves as
the differentiated value our Navy brings to the Joint Force: expeditionary
reach, unmatched mobility, and persistent presence that promotes our vital
interests, protects the global commons, and wins our Nation's wars.
Comprised of People, Platforms, and Payloads, the Fleet spans across all
domains. To our Allies, it serves to reassure. To our Adversaries, it
deters and shapes behavior. To all, it is a signal that freedom of the seas
is non-negotiable and will never be surrendered.
3. As threats evolve and the world grows more complex, we will continue to
build an integrated, all-domain Fleet - capable of conducting sea denial and,
when needed, sea control within the battlespace and at decisive global
chokepoints. This means fusing maritime, cyber, space, and undersea
capabilities, and integrating Joint and combined forces to dominate across
the spectrum of conflict. In short, we will not only design our forces to
fight and win, but to outpace emerging threats through smart and cost-
effective "hedges" to lower the risks faced by our general purposes forces
while increasing survivability and lethality.
4. To that end, building a Fleet to cover every possible threat is too
expensive, unrealistic, and sub-optimized. Instead, I have directed the
development of a new Navy Strategy. Built on the fundamentals from our Navy
Warfighting and Deterrence Concepts, this new effort introduces an enduring
"Hedge Strategy" - designed as an approach that reduces operational risk by
investing in Tailored Offsets and Tailored Forces as a way to amplify and
synergize our general purposes forces to more effectively address low
probability-high consequence contingencies.
5. At its core, our Navy is already a collection of "hedge" forces and
concepts: risk-informed decisions about where to invest and how to fight
while weighing the probability of occurrence in terms of severity and loss.
For example, "Hellscape" is a use case of the generalized Hedge Strategy.
Naval Special Warfare is our hedge against low intensity and irregular
warfare threats. SSBNs and E-6Bs are our hedge that deters strategic attack
with less than eight percent of the Navy. And our Navy-Marine Corps team
itself, has and will continue to serve as the Joint Force's global hedge -
forward deployed, postured, and ready to respond to a variety of threats at a
time and place of our choosing.
6. Tailored Forces are logical groupings of tailored offsets that strengthen
our general purpose forces. The Global Maritime Response Plan (GMRP) and
associated Readiness Conditions (RESCON) provide tailored force packages that
are Combat Surge Ready (CSR) in support of our Combatant Commanders through
tailored certifications to execute specific missions. This is in run today
as we continue to work to achieve our 80 percent CSR goals for all platforms.
7. Further, we will also invest in platforms, sensors, and weapon systems
built with modularity that are scalable and ready for rapid upgrade. We will
accelerate the delivery of integrated, networked capabilities - including
unmanned systems, Artificial Intelligence, and resilient Command and Control
(C2) architectures to ensure decision advantage and guarantee our overmatch
in contested environments around the world.
8. Many of our Navy's toughest challenges cannot be solved in one year or
even across one five-year Future Years Defense Program (FYDP) budgeting
cycle. These types of longer-range problems (e.g., spare parts, critical
munitions, foundry improvements, shipbuilding, and depot repair enhancements)
require time horizons and funding profiles with trajectories that balance
fiscal realities to operational risk.
Therefore, I have directed the development and use of multi-year/multi-FYDP
"funding curves" to ensure sustained investment in the capabilities at the
required capacity our warfighters need.
With clear "commitment to the curve" - I will prioritize the required
sustained investments necessary to solve many of the problems that have
hampered our Navy for decades.
9. The Administrative Chain of Command (ADCON) is comprised of subject
matter experts who form the best force generators in the Navy. Over time we
have drifted into commingling force generation and force employment with
overlapping C2 structures that lack clear accountability. Therefore, we will
review existing C2 structures from an operational task perspective and align
C2 to the risk owner.
Additionally, we will review and align force generation responsibilities to
ensure we are following the "Standard Model" of excellence. This will ensure
that we clearly delineate the role of the appropriate Type Commander (TYCOM)
to the Fleet Commander with clean and accountable hand-offs between force
generation and force employment of their forces.
10. Additionally, to most effectively counter the pacing threat that the
People's Republic of China presents, CTF-54 will be separated from CTF-74
allowing each to hone their focus on specific undersea mission sets,
particularly in our most challenging threat environments. CTF-54 will be
dual-hatted as Commander, Submarine Squadron 21 in Bahrain. To support the
new CTF-54, primary broadcast control authority and submarine operating
authority duties (e.g., waterspace management) will be shifted to Commander,
Submarine Group 8 and CTF-69 in Naples, Italy. To support CTF-74 more fully
across the Western Pacific Area of Responsibility, we will be assessing the
establishment of Task Groups to enable clear and unambiguous C2 lines of
operational accountability.
11. Fleet Improvement Offices (FIOs) are driving a culture of continuous
improvement, intellectually honest assessments, disciplined execution, as
well as improved operational risk management and accountability. We are
codifying this through a new instruction, under the Office of Warfighting
Advantage (OWA), to scale FIO practices across our Foundry, Fleet, and Fight
priorities.
12. Built in the Foundry - Tempered in the Fleet - Forged to Fight.
13. ADM Daryl Caudle, 34th Chief of Naval Operations sends.//
BT
#0001
NNNN
CLASSIFICATION: UNCLASSIFIED/
\end{verbatim}
\end{quote}

\paragraph{Charge of Command (NAVADMIN 225/25).}
\begin{quote}
  \small\begin{verbatim}
CLASSIFICATION: UNCLASSIFIED/
ROUTINE
R 302016Z OCT 25 MID120012171056U
FM CNO WASHINGTON DC
TO NAVADMIN
INFO CNO WASHINGTON DC
BT
UNCLAS
NAVADMIN 225/25
MSGID/NAVADMIN/CNO WASHINGTON DC/CNO/OCT//
SUBJ/CHARGE OF COMMAND - ANNOUNCEMENT TO THE FLEET//
RMKS/1. As the 34th Chief of Naval Operations (CNO), I've released my Charge
of Command for officers selected for command, emphasizing the unique
responsibilities and privileges inherent in leading warfighters at a
strategic inflection point. Commanding Officers (COs) hold the highest
operational and moral accountability for their units, and their leadership
directly shapes combat readiness, culture, and mission success.
2. In the Charge of Command, I characterize command as "an action, not a
position," highlighting the essential attributes of integrity, courage,
humility, and ownership. The Charge emphasizes that the Navy's decisive
advantage is its Sailors, and COs must invest in their people with the same
intensity as operational readiness, fostering trust, accountability, and a
warfighting culture that thrives under pressure.
3. Command inherently involves risk. The Charge directs COs to manage
operational risk deliberately, provide clear Commander's intent, and empower
teams under Mission Command principles.
Disciplined risk-taking, sound judgment, and adaptability are required to
ensure mission success while cultivating initiative and learning throughout
the command.
4. The Charge further stresses the importance of preparation, assessment,
and continuous learning. COs are expected to plan, rehearse, test
assumptions, and innovate faster than potential adversaries. Excellence is
achieved through initiative, critical thinking, disciplined failure as a
learning pathway, and fostering a culture of resilience and innovation.
5. COs are the single accountable officer for their units and are expected
to lead decisively, mentor future leaders, and uphold the Navy's enduring
values and warfighting mindset. COs and their commands must always be ready
to Fight Tonight. As Fleet Admiral Chester Nimitz once said, "When you're in
command, command."
6. Immediate Superior in Command commanders are expected to review the
charge with their subordinate COs, ensuring they understand, internalize, and
apply its principles. This review supports mentorship, reinforces
accountability, and strengthens unit readiness across the fleet.
7. The full Charge of Command letter, including detailed guidance and
reference materials are available at: https://www.navy.mil/
All personnel are encouraged to read, understand, and support the principles
outlined to strengthen their commands and the Navy as a whole.
8. Built in the Foundry - Tempered in the Fleet - Forged to Fight!
9. ADM Daryl Caudle, 34th Chief of Naval Operations sends.//
BT
#0001
NNNN
CLASSIFICATION: UNCLASSIFIED/
\end{verbatim}
\end{quote}
\end{document}
